\chapter{Nebulizer Treatment}

\section*{Overview}
Nebulizer treatment converts liquid medicine into a fine mist that can be inhaled through a mouthpiece or mask. It delivers medicine directly to the airways. The exact approach, setting, and preparation depend on the indication, anatomy, and the patient's health.

\section*{Why It Is Done}
It is commonly used to deliver bronchodilators or other inhaled medicines during asthma attacks, COPD exacerbations, wheezing, and selected respiratory illnesses when a nebulizer is easier to use than an inhaler.

\section*{How It Is Performed}
The prescribed medicine is placed in the nebulizer cup. The patient sits upright, seals the lips around the mouthpiece or wears a mask, and breathes normally until the treatment is complete. The equipment should be cleaned and dried according to instructions.

\section*{Preparation and Precautions}
Use only the prescribed medicine and dose. Tell the clinician about heart rhythm problems, tremor, glaucoma, or medicine allergies. A bronchodilator may cause temporary shakiness, nervousness, palpitations, or a fast heart rate.

\section*{Results and Risks}
Breathing, wheezing, oxygen level, and work of breathing are assessed. Cough, throat irritation, unpleasant taste, or temporary rapid heartbeat may occur. Severe or worsening breathing difficulty requires urgent medical attention rather than repeated home treatments alone.

\section*{When to Seek Medical Care}
Follow the procedure team's specific instructions. Seek urgent medical assessment for severe or worsening pain, uncontrolled bleeding, fever or signs of infection, trouble breathing, chest pain, fainting, new weakness or confusion, inability to urinate, or any rapidly worsening symptom. The appropriate warning signs depend on the procedure and the patient's medical condition.

\section*{References}
\begin{enumerate}
  \item MedlinePlus. Nebulizer treatment. U.S. National Library of Medicine; 2025.
  \item Current Medical Diagnosis and Treatment. McGraw-Hill; 2025.
\end{enumerate}
\clearpage
